\documentclass[a4paper,11pt]{article}
\usepackage[utf8]{inputenc}
\usepackage[T1]{fontenc}
\usepackage[french]{babel}
\usepackage[top=2cm, bottom=2cm, left=2.2cm, right=2.2cm]{geometry}
\usepackage{xcolor}
\usepackage{graphicx}
\usepackage{tcolorbox}
\usepackage{tabularx}
\usepackage{booktabs}
\usepackage{enumitem}
\usepackage{titlesec}
\usepackage{fancyhdr}
\usepackage{listings}
\usepackage{multicol}
\usepackage{caption}
\usepackage{float}
\usepackage{setspace}

% ── Images à placer dans le même dossier que ce .tex ──────────────────
% Renommer vos photos comme suit :
%   img_protova1.jpg   → Protova1 (châssis jaune, Pico visible)
%   img_protova2.jpg   → Protova2 (châssis rouge, Jetson + caméra + antennes)
%   img_g29.jpg        → Volant G29 + pédalier sur support
%   img_ps4.jpg        → Manette PS4 blanche
%   img_rqt.png        → Capture rqt (vue caméra + liste topics)
% ──────────────────────────────────────────────────────────────────────

\definecolor{primary}{HTML}{1A237E}
\definecolor{accent}{HTML}{37474F}
\definecolor{success}{HTML}{2E7D32}
\definecolor{warning}{HTML}{E65100}
\definecolor{lightblue}{HTML}{E3F2FD}
\definecolor{lightgreen}{HTML}{E8F5E9}
\definecolor{lightorange}{HTML}{FFF3E0}
\definecolor{darkgray}{HTML}{424242}
\definecolor{codebg}{HTML}{F5F5F5}

\tcbuselibrary{skins}

\lstset{
  backgroundcolor=\color{codebg},
  basicstyle=\ttfamily\small,
  breaklines=true,
  frame=single,
  rulecolor=\color{gray!40},
  showstringspaces=false,
}

\titleformat{\section}[block]
  {\color{primary}\large\bfseries}{\thesection.}{6pt}{}
  [\vspace{1pt}\color{primary}\hrule\vspace{3pt}]

\titleformat{\subsection}[block]
  {\color{accent}\normalsize\bfseries}{\thesubsection.}{5pt}{}

\pagestyle{fancy}
\fancyhf{}
\fancyhead[L]{\textcolor{primary}{\textbf{Résumé hebdomadaire}} — Semaine du 16 juin 2026}
\fancyhead[R]{\textcolor{darkgray}{\small SEBTI Douae — ALTEN Labs}}
\fancyfoot[C]{\textcolor{darkgray}{\small \thepage}}
\renewcommand{\headrulewidth}{0.3pt}
\setstretch{1.1}

\begin{document}

% ======================== EN-TÊTE ========================
\begin{center}
{\Large\bfseries\color{primary} Résumé Hebdomadaire de Stage}\\[0.3cm]
{\large Semaine du 16 juin 2026}\\[0.15cm]
{\normalsize SEBTI Douae — Stagiaire ALTEN Labs}\\[0.1cm]
{\small\color{darkgray} Encadrant~: Nicolas Baudesson \quad|\quad Tuteur académique~: Samuel Dupont}
\end{center}
\vspace{0.4cm}
\hrule
\vspace{0.5cm}

% ======================== SECTION 1 ========================
\section{Actions réalisées cette semaine}

\subsection{Démonstration de téléopération 5G — Protova2}

La démo de téléopération complète via réseau 5G a été finalisée et validée sur
\textbf{Protova2}. Le système fonctionne avec deux dongles 5G Quectel RM530N-GL
(un côté véhicule, un côté PC opérateur), reliés via un APN privé opérateur
(sous-réseau \texttt{172.16.48.0/28}).

\textbf{Procédure de lancement de la démo~:}
\begin{lstlisting}[language=bash]
# Côté véhicule (Jetson Nano) :
./launch_car.sh     # configure 5G, demarre Zenoh + nodes ROS2

# Côté PC opérateur :
./launch_pc.sh      # configure 5G, demarre Zenoh hub + joy_node + G29
\end{lstlisting}

Le bridge Zenoh s'établit automatiquement entre les deux routeurs via le lien 5G.
Une fois lancé, le pilote peut conduire la voiture RC en temps réel avec le volant
G29, visualiser le flux caméra et superviser les topics ROS2 depuis \texttt{rqt}.

\begin{tcolorbox}[colback=lightgreen, colframe=success, left=4pt, right=4pt,
  top=3pt, bottom=3pt]
\textbf{Performances validées~:} latence ICMP $\approx$ 23--30 ms, débit 5G
disponible 635 Mbps, perte paquets $< 1$\%.
\end{tcolorbox}

\subsection{Téléopération G29 avec retour de force (\textit{force feedback})}

\begin{minipage}[t]{0.62\textwidth}
Le volant Logitech G29 a été intégré avec son \textbf{retour de force} (\textit{force
feedback}) actif pendant la téléopération. Le node \texttt{ros\_g29\_force\_feedback}
reçoit les données du LiDAR (distance frontale) et applique une résistance croissante
sur le volant à mesure que la voiture s'approche d'un obstacle.

\textbf{Rôle du force feedback~:} donner au pilote une sensation haptique de la
situation du véhicule sans qu'il ait besoin de regarder en permanence les
données numériques. Plus l'obstacle est proche, plus le volant résiste —
rendant la conduite à distance plus intuitive et sécurisée, notamment en condition
de latence réseau.

La téléopération PS4 (manette Bluetooth blanche) a également été maintenue en
fonctionnement comme \textbf{alternative sans fil}.
\end{minipage}
\hfill
\begin{minipage}[t]{0.35\textwidth}
\vspace{-4pt}
\begin{center}
\includegraphics[width=\linewidth]{img_g29.jpg}
\captionof{figure}{Volant G29 + pédalier — poste de pilotage ALTEN Labs}
\end{center}
\end{minipage}

\subsection{Visualisation rqt — topics ROS2 actifs pendant la démo}

La figure ci-dessous montre une capture \texttt{rqt} effectuée pendant une session
de téléopération 5G fonctionnelle~: le flux vidéo de la caméra Orbbec Femto Bolt
(RGB-D embarquée sur Protova2) est reçu en direct sur le PC opérateur, et tous les
topics ROS2 sont actifs à travers le bridge Zenoh~: \texttt{/cmd\_vel\_G29},
\texttt{/joy}, \texttt{/camera/color/image\_raw}, \texttt{/ticks},
\texttt{/velocity}, \texttt{/cmd\_vel\_aeb}, etc.

\begin{center}
\includegraphics[width=0.92\textwidth]{img_rqt.png}
\captionof{figure}{Capture \texttt{rqt} — démo 5G en cours~: flux caméra reçu sur PC + liste des topics ROS2 actifs via Zenoh}
\end{center}

\subsection{Présentation des deux prototypes}

\begin{figure}[H]
\centering
\begin{minipage}[b]{0.47\textwidth}
  \includegraphics[width=\linewidth]{img_protova1.jpg}
  \caption{Protova1 — reconstruction depuis zéro (châssis jaune, Raspberry Pi Pico RP2040 visible, phase d'intégration)}
\end{minipage}
\hfill
\begin{minipage}[b]{0.47\textwidth}
  \includegraphics[width=\linewidth]{img_protova2.jpg}
  \caption{Protova2 — plateforme complète (châssis rouge, Jetson Nano, caméra Orbbec, antennes 5G, Seeed Studio)}
\end{minipage}
\end{figure}

\begin{figure}[H]
\centering
\begin{minipage}[b]{0.47\textwidth}
  \includegraphics[width=\linewidth]{img_ps4.jpg}
  \caption{Manette PS4 — téléopération Bluetooth (source secondaire, priorité 60 dans twist\_mux)}
\end{minipage}
\hfill
\begin{minipage}[b]{0.47\textwidth}
  \vspace{10pt}
  \begin{tcolorbox}[colback=lightblue, colframe=primary, left=5pt, right=5pt,
    top=5pt, bottom=5pt]
  \textbf{Arbitrage multi-sources}\\[4pt]
  \small
  Le node \texttt{twist\_mux} arbitre en temps réel entre~:\\[3pt]
  \begin{tabular}{lll}
  AEB sécurité & priorité 255 & (urgence) \\
  PS4 Bluetooth & priorité 60  & \\
  G29 USB       & priorité 30  & \\
  \end{tabular}\\[4pt]
  La priorité maximale de l'AEB garantit que le freinage d'urgence prévaut toujours sur le pilote.
  \end{tcolorbox}
\end{minipage}
\end{figure}


% ======================== SECTION 2 ========================
\section{Perspectives — Semaine du 23 juin 2026}

\subsection{Création de la démo WiFi}

Prochaine étape~: recréer la même démo de téléopération en \textbf{WiFi}, en
parallèle de la démo 5G existante. Les objectifs~:

\begin{enumerate}[noitemsep, leftmargin=1.5cm]
  \item Rédiger un \texttt{launch\_car\_wifi.sh} et \texttt{launch\_pc\_wifi.sh}
        adaptés à la connexion WiFi — ROS2 DDS natif (sans Zenoh), adresses IP du réseau local.
  \item \textbf{Corriger le bug rqt en SSH WiFi}~: actuellement, lancer
        \texttt{rqt} depuis une session SSH ouverte en WiFi génère des erreurs
        (variables \texttt{DISPLAY}/\texttt{XAUTHORITY} non exportées dans le
        contexte SSH). Objectif~: tunnel X11 ou affichage local fonctionnel.
  \item Valider que la démonstration fonctionne de bout en bout en WiFi~:
        G29 téléop + flux caméra + topics actifs dans rqt.
\end{enumerate}

\subsection{Scripts de monitoring WiFi}

Développement de deux nouveaux scripts analogues à ceux de la 5G~:

\begin{center}
\begin{tabularx}{\textwidth}{@{}lXl@{}}
\toprule
\textbf{Script} & \textbf{Rôle} & \textbf{Port} \\
\midrule
\texttt{monitor\_wifi.py} & Monitoring passif~: latence ICMP, débit RX/TX sur interface WiFi, perte paquets, CPU & 8090 \\
\texttt{bench\_wifi\_live.py} & Benchmark actif~: tests iperf3 périodiques UL+DL en WiFi, graphes temps réel & 8091 \\
\bottomrule
\end{tabularx}
\end{center}

\textbf{Note~:} en WiFi, le DDS natif de ROS2 (multicast UDP) fonctionne
directement sur le réseau local — \textbf{Zenoh n'est pas utilisé}. La démo WiFi
repose sur du ROS2 standard, sans bridge supplémentaire. Zenoh est spécifique à la
démo 5G car le multicast DDS ne traverse pas les réseaux cellulaires.

\subsection{Démonstration comparative 5G vs WiFi}

L'objectif final de la semaine est de pouvoir \textbf{passer rapidement de la démo
5G à la démo WiFi devant des clients}, sans manipulation longue. Un script de
bascule rapide sera développé.

\begin{tcolorbox}[colback=lightorange, colframe=warning, left=4pt, right=4pt,
  top=3pt, bottom=3pt]
\textbf{Note sur la comparaison 5G / WiFi~:}
La directrice souhaite mettre en avant \textbf{l'avantage de la 5G sur le WiFi},
notamment en termes de latence et de robustesse du lien. Les scripts
\texttt{monitor\_wifi.py} et \texttt{bench\_wifi\_live.py} permettront d'afficher
côte à côte les métriques des deux réseaux. En conditions réelles de démo
(environnement d'entreprise chargé en interférences WiFi), la 5G sur APN privé
présentera un avantage significatif en latence et stabilité.
\end{tcolorbox}

\subsection{Récapitulatif}

\begin{center}
\begin{tabularx}{\textwidth}{@{}Xcc@{}}
\toprule
\textbf{Tâche} & \textbf{Statut} & \textbf{Délai visé} \\
\midrule
Démo 5G complète (G29 + caméra + force feedback) & \textcolor{success}{\textbf{Terminé}} & Sem. 25 \\
Correction bug rqt en SSH WiFi & En cours & Sem. 26 \\
Démo WiFi (launch files + validation) & Planifié & Sem. 26 \\
\texttt{monitor\_wifi.py} + \texttt{bench\_wifi\_live.py} & Planifié & Sem. 26 \\
Script de bascule 5G ↔ WiFi & Planifié & Sem. 26 \\
Visio de suivi (avant mi-juillet) & À organiser & avant 15 juil. \\
\bottomrule
\end{tabularx}
\end{center}

\vspace{0.6cm}
\hrule
\vspace{0.3cm}
{\small\color{darkgray}
SEBTI Douae — douaesebti.be@gmail.com — ALTEN Labs, Juin 2026
}

\end{document}
